\documentclass[11pt]{article}

\usepackage[margin=1in]{geometry}
\usepackage{enumitem}
\usepackage{amsmath, amssymb}
\usepackage{hyperref}

\setlist[enumerate]{itemsep=0.8em, topsep=0.8em}
\setlist[itemize]{itemsep=0.3em, topsep=0.3em}

\begin{document}

\begin{center}
  {\Large \textbf{CS 421 Examination}}\\[0.25em]
  {\large \textbf{Functional Programming Paradigms}}\\[0.5em]
\end{center}

\noindent\textbf{Instructions:}
\begin{itemize}
  \item Answer all questions.
  \item For Questions 1--4, choose the best option.
  \item For Questions 5--8, write brief but precise answers.
\end{itemize}

\vspace{0.5em}

\begin{enumerate}[label=\textbf{\arabic*.}]

  \item Which property defines a pure function in functional programming?
  \begin{enumerate}[label=(\Alph*)]
    \item It always returns the same output for the same input and has no side effects
    \item It can modify global state but returns consistent results
    \item It uses recursion instead of loops
    \item It operates only on primitive data types
  \end{enumerate}

  \item In the context of lazy evaluation, when are expressions evaluated?
  \begin{enumerate}[label=(\Alph*)]
    \item Immediately when defined
    \item Only when their values are needed
    \item At compile time through optimization
    \item Never, if they are not used in the final output
  \end{enumerate}

  \item What is a monad in functional programming?
  \begin{enumerate}[label=(\Alph*)]
    \item A design pattern for implementing loops
    \item A type class that provides a way to chain operations while abstracting computational context
    \item A data structure for storing immutable collections
    \item A method for parallel computation
  \end{enumerate}

  \item In functional programming, what does the term ``referential transparency'' mean?
  \begin{enumerate}[label=(\Alph*)]
    \item Functions can be replaced by their return values without changing program behavior
    \item Variables can be reassigned multiple times
    \item Functions must explicitly state all dependencies
    \item All data structures are transparent to the compiler
  \end{enumerate}

  \item Explain the concept of higher-order functions and provide examples of how they promote 
  code reusability. How do functions like map, filter, and reduce demonstrate this principle?

  \item Describe the difference between eager and lazy evaluation strategies. What are the 
  advantages and potential pitfalls of lazy evaluation? Provide an example where lazy evaluation 
  enables working with infinite data structures.

  \item Explain the concept of currying and partial application in functional programming. How do 
  these techniques differ, and what practical benefits do they provide in composing functions?

  \item Discuss the concept of immutability in functional programming. Why is immutability considered 
  important, and how does it facilitate reasoning about concurrent programs? What are the trade-offs 
  in terms of performance?

\end{enumerate}

\end{document}
